% 原书分章 PDF 第10—14页：第24.4节及三幅图。
\section{模拟}

在 $Q_0$ 下，基线协变量 $W$ 分解为 $W\equiv(U,V)\in[0,1]\times\{1,2,3\}$，其中 $U$ 和 $V$ 相互独立，分别来自 $[0,1]$ 和 $\{1,2,3\}$ 上的均匀分布。此外，给定 $W$，$Y(0)$ 和 $Y(1)$ 分别从 Beta 分布中抽取；方差固定为 $0.01$，均值分别为 $\bar Q_{0,Y}(0,W)$ 和 $\bar Q_{0,Y}(1,W)$，满足
\[
 \begin{aligned}
 \bar q_{0,Y}(W)&=\bar Q_{0,Y}(1,W)-\bar Q_{0,Y}(0,W)\\
 &\equiv\frac98\left(U^2-\frac52U+\frac23\right)
 +\frac{3\sqrt V}{4\sqrt3}I\left\{U\geq\frac14\left\lfloor\frac{V+3}{3}\right\rfloor\right\},
 \end{aligned}
\]
以及
\[
 \begin{aligned}
 &\bar Q_{0,Y}(1,W)+\bar Q_{0,Y}(0,W)\\
 &\quad\equiv\frac45+\frac1{3\sqrt V}\left(\cos\left(\frac\pi2\frac{4U}V\right)I\{4U\leq V\}\right. \\
 &\hspace{39mm}\left.+\sin\left(\frac\pi2\frac{4U-V}{4-V}\right)I\{4U>V\}-\frac12\right).
 \end{aligned}
\]
条件均值 $\bar Q_{0,Y}(0,\cdot)$、$\bar Q_{0,Y}(1,\cdot)$ 及相应的处理增益函数 $\bar q_{0,Y}$ 见图24.2（左侧两幅）。我们计算以下参数的数值：$\psi_0\approx0.5570$（最优规则 $r_0$ 下的平均奖励）；$\Var_{P_{Q_0,g^b}}\Delta(Q_0,g^b)(O)\approx0.1812^2$（$\Psi$ 在 $P_{Q_0,g^b}$ 处的有效影响曲线在 $P_{Q_0,g^b}$ 下的方差，即在 $Q_0$ 下等概率执行动作 $a=1$ 或 $a=0$ 时的方差）；$\Var_{P_{Q_0,g_0}}\Delta(Q_0,g_0)(O)\approx0.1548^2$（$\Psi$ 在 $P_{Q_0,g_0}$ 处的有效影响曲线在 $P_{Q_0,g_0}$ 下的方差，即在 $Q_0$ 及 $r_0$ 的近似 $g_0$ 下的方差）；以及 $\Var_{P_{Q_0,r_0}}\Delta(Q_0,r_0)(O)\approx0.1512^2$（$\Psi$ 在 $P_{Q_0,r_0}$ 处的有效影响曲线在 $P_{Q_0,r_0}$ 下的方差，即在 $Q_0$ 及 $r_0$ 下的方差）。

设 $p=10\%$、$\xi=1\%$，并选择在 $[-1,1]$ 上由下式刻画的 $G$：
\[
 \begin{aligned}
 G(x)\equiv{}&pI\{x\leq-\xi\}
 +\left(-\frac{1/2-p}{2\xi^3}x^3+\frac{1/2-p}{2\xi/3}x+\frac12\right)I\{-\xi\leq x\leq\xi\}\\
 &+(1-p)I\{x\geq\xi\}.
 \end{aligned}
\]
将 $p$ 降至 $5\%$ 并未显著改变结果（未展示）。工作模型 $\bar{\mathcal Q}_Y$ 由从 $\{0,1\}\times\mathcal W$ 到 $[0,1]$ 的函数 $\bar Q_{Y,\beta}$ 组成：对每个 $a=0,1$ 及 $v\in\{1,2,3\}$，$\logit\bar Q_{Y,\beta}(a,(U,v))$ 是 $1,U,U^2,\ldots,U^5$ 以及 $I\{\frac{j-1}{10}\leq U<\frac j{10}\}$（$1\leq j\leq10$）的线性组合。由此得到的总体参数 $\beta$ 属于 $\mathbb R^{96}$。$\bar Q_{0,Y}$ 和 $\bar q_{0,Y}$ 分别都不属于 $\bar{\mathcal Q}_Y$ 或 $\{\bar q_{Y,\beta}:\bar Q_{Y,\beta}\in\bar{\mathcal Q}_Y\}$。然而，$\expit(\bar q_{Y,0})$ 确实属于后一个工作模型。

在样本量为 $25$ 的倍数且不小于 $200$、不大于 $1000$ 时，执行靶向步骤；实验在样本量达到 $1000$ 时停止。工作模型 $\bar{\mathcal Q}_Y$ 根据负对数似然损失函数 $\ell$，采用惩罚回归进行拟合。通过惩罚，对参数 $\beta$ 的 $\ell_1$ 和 $\ell_2$ 范数施加（数据自适应的）上界，从而寻找一个稀疏最优参数。我们独立地重复第24.1节所述策略 $N=1000$ 次。每次执行靶向步骤时，都构造推论24.2和命题24.3的置信区间，各自的名义覆盖率均设为 $(1-\alpha)=95\%$。模拟研究在 R 中进行（R Development Core Team，2016），使用了 \texttt{tsml.cara.rct} 包（Chambaz，2016）。

\paragraph*{结果。}
图24.1和24.2展示了一次典型实现。图24.2表示 $\bar Q_{0,Y}$、$\bar q_{0,Y}$，及其在最终样本量 $n=1000$ 时的估计量 $\bar Q_{n,Y}$、$\bar q_{n,Y}$。图24.1的上图在每个构造置信区间的样本量 $n$ 处，给出式 (24.10) 中的 $95\%$ 置信区间 $I_n$。根据推论24.2，事件“$\psi_{r_n,0}\in I_n$”的概率，比事件“$\psi_0\in I_n$”的概率更可能接近 $95\%$；这里的意思是，后一性质要求 $\bar q_{n,Y}$ 向 $\bar q_{0,Y}$ 收敛得足够快。尽管如此，在这次实现中，我们看到每个 $I_n$ 都同时包含相应的数据自适应参数 $\psi_{r_n,0}$（叉号）和 $\psi_0$（黑色水平线）。而且，随着 $n$ 增大，$I_n$ 的长度，与样本量 $n$ 处连接同一较深灰度两条曲线的竖直线段长度之间的差异越来越小。这表明，随着 $n$ 增大，$\psi_n^*$ 的方差越来越接近最优方差 $\Var_{P_{Q_0,r_0}}\Delta(Q_0,r_0)(O)$。图24.1的下图在每个构造置信区间的样本量 $n$ 处，给出 $\mathcal R_n$ 的实际值（叉号）及式 (24.11) 中的 $95\%$ 置信区间。我们看到，在这次实现中，遗憾均为正值；这并非事先就有保证。而且，每个置信区间都包含相应的数据自适应参数 $\mathcal R_n$。

\begin{figure}[p]
 \centering
 % 原书真实矢量数据层；中文标签均位于数据区以外。
\begin{tikzpicture}[x=0.47mm,y=-0.47mm]
\node[anchor=north west,inner sep=0pt] at (0,0) {\includegraphics[width=153.2200mm,height=157.4500mm]{figures/ch24_reward_regret_data.pdf}};
\node[font=\fontsize{9.5}{10.5}\selectfont,inner sep=0pt] at (185.1171,156.8511) {靶向步骤的样本量};
\node[font=\fontsize{9.5}{10.5}\selectfont,inner sep=0pt,rotate=90] at (31.2809,72.0720) {[0.95置信区间]};
\node[font=\bfseries\fontsize{10.5}{11.5}\selectfont,inner sep=0pt] at (185.1139,8.1513) {平均奖励};
\node[font=\fontsize{9.5}{10.5}\selectfont,inner sep=0pt] at (185.1217,326.9341) {更新步骤的样本量};
\node[font=\bfseries\fontsize{10.5}{11.5}\selectfont,inner sep=0pt] at (185.1131,178.2341) {遗憾};
\node[font=\fontsize{9.5}{10.5}\selectfont,inner sep=0pt,rotate=90] at (31.2813,242.1546) {[0.95置信区间]};
\end{tikzpicture}

 \caption{最优规则、该规则下平均奖励以及相应遗憾的数据自适应推断示意（另见图24.2）。上图：黑色水平线表示最优规则下平均奖励的值 $\psi_0$。灰色曲线表示映射 $n\mapsto\psi_0\pm\xi_{97.5\%}\sigma_k/\sqrt n$（$k=1,2$），其中 $\sigma_1\approx0.1512$ 为 $\Var_{P_{Q_0,r_0}}\Delta(Q_0,r_0)(O)$ 的平方根（深灰色），$\sigma_2\approx0.1812$ 为 $\Var_{P_{Q_0,g^b}}\Delta(Q_0,g^b)(O)$ 的平方根（浅灰色）。因此，给定样本量 $n$，连接两条深灰色曲线的竖直线段长度，等于基于 $\psi_0$ 的一个正则渐近有效估计量所构造的置信区间长度。叉号表示数据自适应参数 $\psi_{r_n,0}$ 的连续取值。黑点表示 $\psi_n^*$ 的连续取值，以它们为中心的竖直线段表示 $\psi_{r_n,0}$ 的连续 $95\%$ 置信区间；在附加假设下，这些区间也针对 $\psi_0$。下图：叉号表示遗憾 $\mathcal R_n$ 的连续取值。黑点表示 $\psi_n^*-n^{-1}\sum_{i=1}^nY_i$ 的连续取值，竖直线段表示 $\mathcal R_n$ 的连续 $95\%$ 置信区间。}
 \label{fig:ch24-reward-regret}
\end{figure}

\begin{figure}[p]
 \centering
 % 原书真实矢量数据层；中文标签均位于数据区以外。
\begin{tikzpicture}[x=0.46mm,y=-0.46mm]
\node[anchor=north west,inner sep=0pt] at (0,0) {\includegraphics[width=153.6400mm,height=155.4800mm]{figures/ch24_conditional_means_data.pdf}};
\node[font=\bfseries\fontsize{10.5}{11.5}\selectfont,inner sep=0pt] at (97.2533,8.0453) {真值};
\node[font=\bfseries\fontsize{10.5}{11.5}\selectfont,inner sep=0pt] at (271.0629,8.0453) {估计值};
\end{tikzpicture}

 \caption{通过条件均值 $\bar Q_{0,Y}$、处理增益函数 $\bar q_{0,Y}$ 及其估计量，展示最优规则、该规则下平均奖励以及相应遗憾的数据自适应推断（另见图24.1）。左上图：实线表示 $U\mapsto\bar Q_{0,Y}(1,(U,v))$，分别取 $v=1$（黑色，在 $1$ 处最低）、$v=2$（深灰色，在 $1$ 处居中）和 $v=3$（浅灰色，在 $1$ 处最高）。虚线表示 $U\mapsto\bar Q_{0,Y}(0,(U,v))$，分别取 $v=1$（黑色，在 $1$ 处最高）、$v=2$（深灰色，在 $1$ 处居中）和 $v=3$（浅灰色，在 $1$ 处最低）。左下图：曲线表示 $U\mapsto\bar q_{0,Y}(U,v)$，分别取 $v=1$（黑色，在 $1$ 处最低）、$v=2$（深灰色，在 $1$ 处居中）和 $v=3$（浅灰色，在 $1$ 处最高）。右侧两图：与左侧两图相对应，将 $\bar Q_{0,Y}$ 和 $\bar q_{0,Y}$ 换为最终样本量 $n=1000$ 时的 $\bar Q_{n,Y}$ 和 $\bar q_{n,Y}$。}
 \label{fig:ch24-conditional-means}
\end{figure}

我们可以评价这些 $95\%$ 置信区间是否达到了名义覆盖率 $95\%$。为此，进行二项检验。由构造可知，覆盖 $\psi_{r_n,0}$ 的置信区间的经验个数，是来自参数为 $(N,\pi)$ 的二项分布的随机变量。我们检验原假设“$\pi\geq95\%$”，对立假设为单侧的“$\pi<95\%$”。较大的 $p$ 值被解释为：缺乏支持该置信区间未达到其名义覆盖率的经验证据。对 $\psi_0$ 和 $\mathcal R_n$ 也作相应的同样处理。

我们并不报告 $3\times33=99$ 个经验覆盖比例及相应 $p$ 值，而仅绘制评价 $\psi_{r_n,0}$ 与 $\psi_0$ 覆盖率的检验 $p$ 值的对数，见图24.3。总体而言，灰色曲线高于黑色曲线，表明 $\psi_{r_n,0}$ 的经验覆盖率往往较高（范围为 $0.917$ 至 $0.955$，平均为 $0.940$），而 $\psi_0$ 的经验覆盖率较低（范围为 $0.919$ 至 $0.946$，平均为 $0.937$）。正如本节第一段所论述的，这并不令人意外。此外，大多数 $p$ 值大于 $5\%$（上方灰色水平线），而更多 $p$ 值大于 Bonferroni 校正后的阈值 $5/33\%$。最小的 $p$ 值对应于样本量 $n=200$ 和 $n=225$，此时推断所依据的信息很少。$\mathcal R_n$ 的覆盖率远高于名义覆盖率 $95\%$，介于 $0.951$ 与 $0.990$ 之间，平均为 $0.997$。这同样不意外，因为其置信区间按构造是保守的。

\begin{figure}[htbp]
 \centering
 % 原书真实矢量数据层；中文标签均位于数据区以外。
\begin{tikzpicture}[x=0.46mm,y=-0.46mm]
\node[anchor=north west,inner sep=0pt] at (0,0) {\includegraphics[width=152.2600mm,height=78.2000mm]{figures/ch24_p_values_data.pdf}};
\node[font=\fontsize{9.5}{10.5}\selectfont,inner sep=0pt] at (183.9776,162.9361) {更新步骤的样本量};
\node[font=\fontsize{9.5}{10.5}\selectfont,inner sep=0pt,rotate=90] at (9.6152,73.3002) {$\log_{10}(p\text{值})$};
\node[font=\fontsize{9.5}{10.5}\selectfont,inner sep=0pt,rotate=90] at (18.6404,73.2977) {经验覆盖率的二项检验};
\end{tikzpicture}

 \caption{置信区间覆盖率的经验评价。曲线表示二项检验 $p$ 值的对数，原假设为覆盖率充分，对立假设为覆盖率不足。较大的 $p$ 值被解释为：缺乏支持相应置信区间未达到名义覆盖率 $95\%$ 的经验证据。黑色曲线对应 $\psi_{r_n,0}$ 的置信区间，灰色曲线对应 $\psi_0$ 的置信区间。灰色水平线表示 $5\%$ 的阈值（上方），以及 Bonferonni 校正后的阈值 $5/33\%$（下方）。}
 \label{fig:ch24-coverage}
\end{figure}
\FloatBarrier
